\documentclass[10pt,a4paper]{amsart}
\usepackage[margin=19mm]{geometry}
\usepackage{amsmath,amssymb,mathtools}
\usepackage[scr=rsfs,cal=euler]{mathalfa}
\usepackage{xfrac}
\usepackage[unicode,hidelinks,bookmarks=false]{hyperref}
\newcommand{\ie}{i.e.\ }
\newcommand{\cf}{cf.\ }
\newcommand{\resp}{respectively\ }
\newcommand{\Subsection}{\subsection}
\providecommand{\nobreakdash}{\nobreak\hbox{-}}
\setlength{\emergencystretch}{2em}
\multlinegap0pt
\usepackage{amssymb}
\usepackage{tikz}
\newtheorem{prop}{Proposition}
\renewcommand{\H}{\mathcal{H}}
\newcommand{\lPi}{\rotatebox[origin=c]{180}{$\mathsf{\Pi}$}} %{\amalg}%{\underline{\Pi}}
\newcommand{\uPi}{\mathsf{\Pi}}%{\Pi}%{\overline{\Pi}}
\title{Possibilistic inferential models, induction, and the rule of succession}
\author{Ryan Martin, Shih-Ni Prim, Max Raner, Jonathan P. Williams}
\date{Source: arXiv:2608.11935v2 (CC BY 4.0)}
\begin{document}
\maketitle

\noindent\emph{Source and licence.} Possibilistic inferential models, induction, and the rule of succession. Version arXiv:2608.11935v2 (CC BY 4.0), distributed by arXiv under the Creative Commons Attribution 4.0 International licence, http://creativecommons.org/licenses/by/4.0/, as recorded in the arXiv licence metadata for this exact version and shown on the abstract page https://arxiv.org/abs/2608.11935v2. Full licence notice: \texttt{licenses/arxiv-2608.11935-CC-BY-4.0.txt}.\par\smallskip

\noindent\textit{Editorial scope.} Counted unit: the article's prop prop:monotone (source0.tex lines 345--348). The article's complete original proof of it is delivered with it (source0.tex lines 350--357, one \texttt{proof} environment of 8 lines). No mathematics is written, repaired or re-derived: the statement and the proof are the article's own lines, in the article's own notation, and no retained line is substituted or re-flowed.  No further source-local context is needed: every cross-reference of every retained line resolves inside the two delivered spans.  Nothing was omitted from any retained span, so the package carries no omission record.  No retained line contains a citation, so the wrapper carries no bibliography and no bibliography entry.  No retained line contains a figure environment, an image-inclusion command or drawing code, so no image is delivered and the package declares no asset dependency.  Editorial formatting only, in addition: an AMS \texttt{amsart} preamble that restates the article's own declarations and macros used by the retained lines, namely the environments \texttt{prop} on separate counters, together with the standard packages those lines need.  The retained environments print in this document as Proposition 1 in order of appearance, whereas the article prints the same items with the numbering of its own full document.  The complete native source of the article as distributed in the arXiv e-print for this exact version is bundled unchanged as \texttt{sources/207459/source0.tex}.
\begin{prop}
\label{prop:monotone}
Given $h$, $\lPi_x(h \lor h') > \lPi_x(h)$ if and only if the best explanation $\eta \in \H_\text{\rm elem}$ of $x$ that entails $\neg h$ does not entail $\neg h \land \neg h'$.
\end{prop}

\begin{proof}
It is easy to see that 
\begin{align*}
\lPi_x(h \lor h') > \lPi_x(h) & \iff \uPi_x(\neg h \land \neg h') < \uPi_x(\neg h) \\
& \iff \max_{\eta \in \H_\text{elem}: \eta \models \neg h \land \neg h'} \pi_x(\eta) < \max_{\eta \in \H_\text{elem}: \eta \models \neg h} \pi_x(\eta). 
\end{align*}
The latter inequality holds if and only if the maximum on the right is attained at $\eta$ with $\eta \models \neg h$ and that this same $\eta$ has $\eta \not\models \neg h \land \neg h'$ or, in other words, if the best explanation $\eta \in \H_\text{\rm elem}$ of $x$ that entails $\neg h$ does not entail $\neg h \land \neg h'$.
\end{proof}

\end{document}
